\textbf{1. Kiến thức cần nhớ}
\begin{itemize}
\item \textbf{Đường trung bình của tam giác} là đoạn thẳng nối trung điểm hai cạnh của tam giác. Mỗi tam giác có ba đường trung bình.
\item \textbf{Định lí}: Đường trung bình của tam giác song song với cạnh thứ ba và bằng nửa cạnh đó.\\
$M$, $N$ lần lượt là trung điểm của $AB$, $AC$ $\Rightarrow$ $\boxed{MN\parallel BC,\ MN=\dfrac12BC}$.
\item \textbf{Nhận xét} (suy từ định lí Thalès): đường thẳng đi qua trung điểm một cạnh của tam giác và song song với cạnh thứ hai thì đi qua trung điểm cạnh thứ ba.
\end{itemize}

\textbf{2. Phương pháp giải các dạng}

%% DANG 8H4?16-1
\textbf{Dạng 1. Tính độ dài nhờ đường trung bình}
\begin{enumerate}
\item Chỉ ra hai trung điểm của hai cạnh $\Rightarrow$ đoạn nối chúng là đường trung bình.
\item Áp dụng $MN=\dfrac12BC$ (hoặc $BC=2MN$) để tính độ dài.
\end{enumerate}
Lưu ý: đường trung bình bằng nửa cạnh \textit{thứ ba} (cạnh không chứa hai trung điểm).

%% DANG 8H4?16-2
\textbf{Dạng 2. Chứng minh trung điểm, song song nhờ đường trung bình}
\begin{enumerate}
\item Chứng minh song song: chỉ ra đoạn thẳng là đường trung bình của một tam giác $\Rightarrow$ song song với cạnh thứ ba.
\item Chứng minh trung điểm: dùng nhận xét — qua trung điểm một cạnh, kẻ song song với cạnh thứ hai thì cắt cạnh thứ ba tại trung điểm.
\item Bài tổng hợp: có thể dùng đường trung bình của nhiều tam giác (ví dụ tứ giác có các trung điểm $\Rightarrow$ hình bình hành).
\end{enumerate}
